% Fragmento público generado desde propietarios íntegros.
% Residencia: 52.4.1.
% Fuente: ../incoming/radion_monografia_20260827/manuscrito/sections/05_geometria_accion.tex:1-217
\noindent La geometría espectral determina la forma y la ley areal determina la acción transportada por la fase.\par\medskip

\subsubsection{Equivalencia entre las cartas espectral y geométrica}

\paragraph{Operador espectral positivo y autovalores conjugados}

El par angular produce el operador bicapa
\begin{equation}
B_1=AI+C^*R,\qquad R^2=I,\qquad
P_\pm=\frac{I\pm R}{2}.
\label{eq:B1}
\end{equation}
Sus autovalores son \(A\pm C^*\). Como \(0<C^*<A\), ambos son positivos.
En la carta \(\kappa=\pi/180\), la elipse espectral tiene semiejes
\begin{equation}
a_1^2=\kappa(A+C^*),\qquad
b_1^2=\kappa(A-C^*).
\label{eq:semiejes-espectrales}
\end{equation}
La descomposición es reversible:
\begin{equation}
A=\frac{a_1^2+b_1^2}{2\kappa},\qquad
C^*=\frac{a_1^2-b_1^2}{2\kappa}.
\label{eq:reconstruccion-AC}
\end{equation}
La elipse no se dibuja para ilustrar un número. Reconstruye exactamente el
par que la generó.

Definamos
\begin{equation}
\rho_1=\sqrt{A^2-C^{*2}},\qquad
\eta=\operatorname{artanh}\frac{C^*}{A}
=\frac12\log\frac{A+C^*}{A-C^*}.
\label{eq:eta}
\end{equation}
Entonces
\[
A\pm C^*=\rho_1e^{\pm\eta},
\]
y
\begin{equation}
a_1=\sqrt{\kappa\rho_1}\,e^{\eta/2},
\qquad
b_1=\sqrt{\kappa\rho_1}\,e^{-\eta/2}.
\label{eq:semiejes-eta}
\end{equation}
El producto conserva escala; el cociente conserva forma:
\[
a_1b_1=\kappa\rho_1,
\qquad
\frac{b_1}{a_1}=e^{-\eta}
=\sqrt{\frac{A-C^*}{A+C^*}}.
\]

\paragraph{Invariantes focales del operador positivo}

La semidistancia focal es
\begin{equation}
c_1^2=a_1^2-b_1^2=2\kappa C^*,
\qquad
F_\pm=(\pm c_1,0).
\label{eq:focos}
\end{equation}
La distancia total vale
\begin{equation}
d_1=2c_1=2\sqrt{2\kappa C^*}
=0.544545509325626623406299779866023\ldots.
\label{eq:d1}
\end{equation}
El número aislado no es la tesis. La identidad \(c_1^2=2\kappa C^*\) dice que
\(C^*\) es literalmente la escisión focal cuadrática en la carta
distinguida. Los focos son los extremos conjugados de la apertura espectral;
no son, por sí solos, los polos eléctrico y magnético. Esa lectura requiere
el transductor constitutivo.

La excentricidad es
\begin{equation}
e_f=\frac{c_1}{a_1},\qquad
e_f^2=1-e^{-2\eta}=\frac{2C^*}{A+C^*}.
\label{eq:excentricidad}
\end{equation}
Numéricamente,
\[
\begin{aligned}
\eta&=0.299689683921630799684926562863927\ldots,\\
e_f&=0.671451895550191986916277055864332\ldots.
\end{aligned}
\]

\paragraph{Bloque APP de referencia y transporte normalizado}

El bloque central
\[
B_0=7I+2R=9P_++5P_-
\]
posee
\[
\rho_0=\sqrt{45},\qquad
\eta_0=\frac12\log\frac95,\qquad
e_0=\frac23.
\]
Su distancia focal es
\[
d_0=4\sqrt\kappa
=0.528443639680801468446003835349358\ldots.
\]
La comparación exacta es
\begin{equation}
\boxed{\left(\frac{d_1}{d_0}\right)^2=\frac{C^*}{2}},
\qquad
d_1=d_0\sqrt{\frac{C^*}{2}}.
\label{eq:invariante-focal-relativo}
\end{equation}
Ésta es la genealogía defendible de \(0.544545\ldots\): una escala focal APP
de referencia multiplicada por un invariante relativo. Su cercanía visual a
\(0.54\) no es una identidad con el defecto APP \(54\).

El transporte desde \(B_0\) se separa en escala y anisotropía:
\begin{equation}
T_{\mathrm{rel}}=\frac{\rho_1}{\sqrt{45}}
\exp\bigl((\eta-\eta_0)R\bigr).
\label{eq:transporte-rel}
\end{equation}
Sobre las hojas,
\[
t_+=\frac{A+C^*}{9}=1.0467878786947163\ldots,
\qquad
t_-=\frac{A-C^*}{5}=1.0347228460630311\ldots.
\]
La deformación no es un solo decimal. Su parte isotrópica es
\(\sqrt{t_+t_-}\); su parte hiperbólica es
\(\tfrac12\log(t_+/t_-)\).

\subsubsection{Elipse positiva asociada al par angular}

\paragraph{Construcción mediante semiejes conjugados}

La escala espectral anterior carece de unidades de acción. Mecánica
Dimensional publica la misma forma sobre la recta interna de acción:
\begin{equation}
a_S=\sqrt{2\hret e^\eta},
\qquad
b_S=\sqrt{2\hret e^{-\eta}}.
\label{eq:semiejes-accion}
\end{equation}
Inmediatamente,
\begin{equation}
a_Sb_S=2\hret,\qquad
\frac{a_S}{b_S}=e^\eta.
\label{eq:producto-forma}
\end{equation}
El producto fija la acción; el cociente fija la anisotropía. La elipse se
parametriza por
\begin{equation}
Q(\theta)=a_S\cos\theta,\qquad
P(\theta)=b_S\sin\theta.
\label{eq:elipse-param}
\end{equation}

Las elipses espectral y de acción son homotéticas:
\begin{equation}
\lambda_{\mathrm{act}}^2=\frac{2\hret}{\kappa\rho_1},
\qquad
(a_S,b_S,c_S)=\lambda_{\mathrm{act}}(a_1,b_1,c_1).
\label{eq:homotecia}
\end{equation}
Comparten la forma \(e^{-\eta}\), no la dimensión física.

\paragraph{Ley areal de la acción por radión}

La forma simpléctica canónica es \(\omega=\diff Q\wedge\diff P\). El área
orientada barrida desde \(0\) hasta \(\theta\) es
\begin{align}
\mathcal A(\theta)
&=\frac12\int_0^\theta
\bigl(Q(t)P'(t)-P(t)Q'(t)\bigr)\,\diff t\\
&=\frac12a_Sb_S\int_0^\theta
(\cos^2t+\sin^2t)\,\diff t.
\end{align}
Usando \eqref{eq:producto-forma}, obtenemos el resultado nuclear.

\begin{theorem}[Acción por radión]
Para la elipse \eqref{eq:elipse-param},
\begin{equation}
\boxed{\mathcal A(\theta)=\hret\theta.}
\label{eq:area-theorem}
\end{equation}
En particular,
\begin{equation}
\boxed{\mathcal A(1)=\hret},
\qquad
\boxed{\mathcal A(2\pi)=2\pi\hret=\hfull}.
\label{eq:area-radion-vuelta}
\end{equation}
\end{theorem}

Esta igualdad permite definir el radión sin apelar primero a la longitud de
arco.

\begin{definition}[Radión orientado]
\begin{equation}
\mathfrak r_\sigma=(\theta=1,\sigma),
\qquad \sigma\in\{+1,-1\},
\qquad
\mathcal A(\mathfrak r_\sigma)=\sigma\hret.
\label{eq:radion-oriented}
\end{equation}
\end{definition}

El radián es la proyección que olvida \(\sigma\), la hoja, la memoria y el
área. El radión conserva esa información. De ahí el nombre \emph{rad--ión}:
la orientación precede a la carga física, y la carga aparece más tarde como
\[
Q(x)=n_Q(x)e_{\mathrm{int}},\qquad
e_{\mathrm{int}}=\sqrt{\frac{2\alphah\hfull}{Z_{\mathrm{int}}}}.
\]

